\documentclass{article}

\usepackage{booktabs}
\usepackage{tabularx}
\usepackage{enumitem}
\usepackage{float}

\setlength{\emergencystretch}{3em}

\title{Development Plan\\\progname}

\author{\authname}

\date{}

\input{../Comments.text}
\input{../Common.text}

\begin{document}

\maketitle

\begin{table}[hp]
\caption{Revision History} \label{TblRevisionHistory}
\begin{tabularx}{\textwidth}{llX}
\toprule
\textbf{Date} & \textbf{Developer(s)} & \textbf{Change}\\
\midrule
Sept 22, 2026 & William Pol & Initial draft for TA meeting\\
Sept 27, 2026 & William Pol & Revised against the Development Plan checklist
and rubric; added PoC risks, workflow details, role rotation plan, and
course schedule\\
Sept 28, 2026 & William Pol & Addressed pull request review; completed
the team charter and reflection\\
Nov 2026 & William Pol & Updated for cross-model KV cache translation:
terms, technical areas, schedule, PoC plan and technology; meetings after
TA and supervisor meetings; reflection questions restored\\
\bottomrule
\end{tabularx}
\end{table}

\newpage{}

\section*{About This Document}

This document records how Team 11 will build \progname{}: how we meet and
communicate, who is responsible for what, how code and documents move
through git, when major pieces of work happen, how the Proof of Concept
(PoC) demonstration will test our biggest risk, and which tools we expect
to use. The problem and goals themselves are described in the Problem
Statement and Goals document. This plan is written for two audiences: the
current team, as the reference for how we agreed to work, and viewers of
the document, such as the teaching assistant (TA), the instructor, and
peer reviewers, who
should be able to read it and understand how the project is run without
asking the team directly. Decisions recorded here can change; changes are
made through a pull request that updates this document.

Table~\ref{TblTerms} defines the terms and acronyms used in this
document. Each acronym is also spelled out where it is first used.

\begin{table}[H]
\caption{Terms and acronyms} \label{TblTerms}
\begin{tabularx}{\textwidth}{lX}
\toprule
\textbf{Term} & \textbf{Meaning}\\
\midrule
LLM & Large language model.\\
Semantic routing & Choosing which model answers a request based on what
the request asks for.\\
Prefill & The first phase of LLM inference, in which the model processes
the whole prompt at once.\\
Decode & The phase after prefill, in which the model generates output one
token at a time.\\
KV cache & The keys and values computed for each prompt token during
prefill, stored so they do not have to be recomputed during decode.\\
TTFT & Time to first token: the time from receiving a request to producing
the first output token.\\
Source / target model & On a model switch, the model that built the
existing KV cache, and the model that answers next.\\
Translator & A small trained model that turns the source model's KV cache
into a KV cache the target model can use.\\
Paired KV data & KV caches that both models produce for the same text,
used to train and evaluate the translator.\\
Fallback & Running the target model's normal prefill when a translated
cache is unavailable or rejected.\\
vLLM & An open-source LLM serving engine that Tandem will integrate with
after the PoC.\\
\bottomrule
\end{tabularx}
\end{table}

\section{Confidential Information?}

This project does not involve any confidential information from industry.
The team's supervisor, Manoj Govindassamy, is from industry (Cohere), but
the project uses no confidential information from Cohere, so there is no
non-disclosure agreement (NDA) or other confidentiality agreement in place.

\section{IP to Protect}

There is no intellectual property (IP) to protect. \progname{} is an
open-source systems project that builds on existing open-source large
language models (LLMs) and inference libraries, with no industry sponsor, so no IP
agreement or
Intellectual Property Guide Acknowledgement is required.

\section{Copyright License}

The team is adopting the MIT License. The license text, with the team
members as copyright holders, is in the \texttt{LICENSE} file at the root
of the \href{https://github.com/jjustin-k/Tandem}{project repository}.

\section{Team Meeting Plan}

\paragraph{When and where.} The team meets immediately after each TA
meeting and each supervisor meeting, to catch up on work and work through
blockers. A meeting after a TA meeting is in person on campus, with the
members who attended the TA meeting. Supervisor meetings, and the team
meetings after them, are online through Google Meet. Physical meeting locations are shared in the team
chat, not in this public repository. When a lecture or tutorial slot has
nothing scheduled, the team uses it as an extra working session. In the
week before a deliverable is due, the team adds a short check-in.

\paragraph{Agenda and minutes.} A team meeting has no issue of its own,
since it is part of the TA or supervisor meeting it follows. That meeting's
GitHub issue, created from the \texttt{TA Meeting} or \texttt{Supervisor
Meeting} issue template, holds the agenda and the attendance. During the
meeting Granola records notes and a summary automatically; the chair adds
Granola's notes to the issue as a comment, with the decisions and action
items from both the TA or supervisor meeting and the team meeting. Every
action item becomes its own issue, assigned to its owner, and the meeting
issue is closed once those issues exist.

\paragraph{Chair.} William Pol chairs the team meetings.

\paragraph{Supervisor meetings.} The team's supervisor is Manoj
Govindassamy (Cohere). The team meets with him once a month on Google
Meet, tracked with the \texttt{Supervisor Meeting} issue template.

\section{Team Communication Plan}

\begin{itemize}
  \item \textbf{Microsoft Teams} is for day-to-day discussion and quick
  questions; online meetings use Google Meet. Members are expected to reply to
  messages that mention them within 24 hours on weekdays. Anything urgent
  (for example, a broken build close to a deadline) is posted with an
  @-mention of the people needed.
  \item \textbf{GitHub issues and pull requests} are the record of work.
  Any decision made in chat or in a meeting that affects the code or the
  documents is written into the relevant issue or PR, so the reason
  behind it can be found later. If it is not in GitHub, it was not
  decided.
  \item \textbf{The TA and instructor} are contacted by the team liaison,
  Justin Kwinecki, with the rest of the team copied. Course discussions on
  the capstone GitLab repository are handled the same way.
  \item \textbf{Files} (documents, code, data, results) live in the
  repository, not in chat, so there is one current copy.
\end{itemize}

\section{Team Member Roles}
\label{sec:roles}

Roles are split into project roles, which keep the team running, and
technical areas, which cover the expertise the project needs. No one owns a
single narrow role for the whole year.

\subsection{Project Roles}

Table~\ref{TblProjectRoles} lists the roles that keep the team running.

\begin{table}[H]
\caption{Project roles} \label{TblProjectRoles}
\begin{tabularx}{\textwidth}{lX}
\toprule
\textbf{Role} & \textbf{Responsibility and assignment}\\
\midrule
Liaison & Point of contact with the instructor and TA. Justin Kwinecki,
for the whole year.\\
Meeting chair & Writes the agenda, runs the meeting, ticks off attendance,
and turns Granola's notes into decisions and action items. William Pol.\\
Reviewer & Reviews a pull request before it is merged. Anyone except the
author; the author requests a reviewer from outside their technical area
at least once a week, to spread knowledge.\\
Quality assurance (QA) & For each deliverable, one member who did not
write most of it checks it against the matching checklist in
\texttt{docs/Checklists/} and the course rubric, and files issues for
anything missing. Rotates per deliverable.\\
Deliverable lead & For each document or demo, one member coordinates who
writes which part and makes sure it is submitted on time. Rotates per
deliverable, so every member leads at least one each term.\\
\bottomrule
\end{tabularx}
\end{table}

\subsection{Technical Areas}

Table~\ref{TblAreas} lists the three broad areas of technical work.
Documentation is shared by everyone and coordinated by the deliverable
lead.

\begin{table}[H]
\caption{Technical areas} \label{TblAreas}
\begin{tabularx}{\textwidth}{lX}
\toprule
\textbf{Area} & \textbf{What it covers}\\
\midrule
Translator and data & Generating paired KV data from both models,
training the translator, and evaluating it against native prefill.\\
Router and runner & The prompt classifier and routing policy, the
inference runner, and handing a translated cache to the target model.\\
Serving integration and measurement & Integration with vLLM, benchmarks
and experiments, access to the department GPU clusters, the Python
environment, and continuous integration (CI).\\
\bottomrule
\end{tabularx}
\end{table}

\paragraph{How technical areas are assigned.} Before the PoC, every
member worked on the PoC so that everyone learned how the models, their KV
caches and the translator fit together. The areas are now owned as
follows:
\begin{itemize}[noitemsep]
  \item Translator and data: Zayed Umer (training and evaluation) and
  Musab Mondal (the runtime adapter that applies the translator).
  \item Router and runner: Justin Kwinecki (the runner and the demo) and
  Manav Patel (choosing the translator in the runner).
  \item Documentation and design: William Pol, with the deliverable lead
  for each document.
\end{itemize}
Serving integration and measurement starts after the PoC; its owner is
chosen at the first team meeting after the PoC demo. Every area has at
least two people who understand it.

\paragraph{Changing roles.} Roles are reviewed at the first team meeting
after each major demo (PoC, Rev 0). Any member can also ask to change areas
at a team meeting. When an area changes hands, the outgoing owner writes
a short handover issue (current state, open problems, where the code is)
and pairs with the new owner for one week.

\section{Workflow Plan}
\label{sec:workflow}

This section is written to become the project's \texttt{CONTRIBUTING.md}.

\subsection{Branches and Commits}

\begin{itemize}
  \item \texttt{main} is protected: nothing is pushed to it directly, and
  every change arrives through a pull request (PR).
  \item Work happens on a short-lived branch named
  \texttt{type/short-description}. The type is one of \texttt{feat},
  \texttt{fix}, \texttt{docs}, \texttt{poc}, \texttt{bench}, or
  \texttt{chore}.
  \item Commits are small and focused, with a message in the imperative
  mood that explains why the change was made, not just what changed.
\end{itemize}

\subsection{Pull Requests}

\begin{itemize}
  \item Every PR links the issue it resolves (``Closes \#12''), describes
  what changed, and says how it was tested.
  \item A PR needs one approving review from another team member and
  passing CI before it is merged. Work in progress is opened as a draft PR
  so others can see it early.
  \item Reviewers check correctness and readability, and for documents,
  the matching checklist in \texttt{docs/Checklists/}.
  \item PRs are merged with ``squash and merge'' so each PR becomes one
  commit on \texttt{main}, and the branch is deleted afterwards.
\end{itemize}

\subsection{Tags}

When a deliverable is submitted, the commit on \texttt{main} that was
submitted is marked with an annotated tag (for example, \texttt{PS-DevPlan},
\texttt{SRS-HA}, \texttt{PoC}, \texttt{Rev0}, \texttt{Final}). This lets
the TA, the instructor, and peer reviewers see exactly what was graded, and
lets the team compare later work against it.

\subsection{Issues}

\begin{itemize}
  \item Every piece of work, bug, research question, and TA or supervisor
  meeting is an issue. The repository already has issue templates for team meetings, TA
  meetings, supervisor meetings, lectures, and peer reviews.
  \item Each issue gets one \emph{type} label and, for code, one
  \emph{area} label:
  \begin{itemize}[noitemsep]
    \item Type: \texttt{bug}, \texttt{enhancement}, \texttt{documentation},
    \texttt{research}, \texttt{meeting}, \texttt{review}, \texttt{question}.
    \item Area: one per technical area: translator and data, router and
    runner, and serving integration and measurement.
  \end{itemize}
  \item Each issue has one assignee and belongs to a milestone. There is
  one milestone per course deliverable, with the deliverable's due date.
  \item Peer review issues from other teams (label \texttt{review}) are
  answered within a week and closed with a link to the commit that
  addresses them, or a short explanation if no change is made.
\end{itemize}

\subsection{Continuous Integration}

\begin{itemize}
  \item \textbf{Documents (already running):} a GitHub Actions workflow
  (\path{.github/workflows/latex-pages.yml}) rebuilds every changed
  document, commits the PDFs to \texttt{pdfs/}, and publishes them to the
  project's GitHub Pages site.
  \item \textbf{Code (added with the first code PR):} a second workflow
  runs the linter and the unit tests on every PR, including the existing
  router, classifier and runner tests in \texttt{test/}, on the central
  processing unit (CPU). This workflow is not set up yet.
  \item \textbf{GPU work:} GitHub's CI machines have no GPU, so GPU tests
  and benchmarks are run by hand on the department clusters, and their
  results are committed to the repository with the command used to produce
  them.
\end{itemize}

\section{Project Decomposition and Scheduling}

The team uses a GitHub Project board,
\href{https://github.com/users/jjustin-k/projects/11}{Tandem Kanban}, with
the columns Backlog, Ready, In Progress, In Review, and Done. Every issue is
on the board, and the board is reviewed at each team meeting.

Table~\ref{TblSchedule} combines the course deadlines with our technical
milestones. Dates for Term 1 are from the course's weekly announcements;
later dates follow the course schedule and will be updated as they are
announced.

\begin{table}[ht]
\caption{Schedule} \label{TblSchedule}
\begin{tabularx}{\textwidth}{lX}
\toprule
\textbf{When} & \textbf{Course deliverable (and our technical work)}\\
\midrule
Sept 28 & Problem Statement, PoC Plan, Development Plan, Team Charter\\
Oct 1 & Peer review of another team's Problem Statement and Plans\\
Oct 19 & Software Requirements Specification (SRS) and Hazard Analysis\\
Oct 22 & Peer review of another team's SRS and Hazard Analysis\\
Week 7 & Verification and Validation (V\&V) Plan; PoC schedule requests\\
Week 8 & Design Document Rev $-1$ and Fall custom document (machine
learning (ML) model report); PoC productivity report\\
Mid--late Nov & \textbf{PoC demonstration}\\
\midrule
Sept--Nov & \emph{PoC work, each item due by Nov 24, the day before the
demo: paired KV dataset and the SmolLM to Qwen translator (Zayed Umer);
inference runner and demo (Justin Kwinecki); translator runtime adapter
(Musab Mondal); translator choice in the runner (Manav Patel); end-to-end
check of the translated handoff (owner to be assigned)}\\
Dec--Jan & \emph{Qwen to SmolLM translator; vLLM integration}\\
\midrule
Jan (Week 16) & Design Document Rev 0; module implementation\\
Feb (Weeks 18--19) & \textbf{Rev 0 demonstration}\\
\midrule
Feb--Mar & \emph{Serving benchmarks against native prefill on the
department GPU clusters}\\
\midrule
Mar & V\&V Report and Winter custom document (Performance Report)\\
Late Mar & \textbf{Final demonstration}\\
Apr & Final documentation, Capstone EXPO poster and video, EXPO\\
\bottomrule
\end{tabularx}
\end{table}

\section{Proof of Concept Demonstration Plan}
\label{sec:poc}

The PoC demonstration tests whether one model's KV cache can be
translated for another model well enough, and fast enough, to replace the
target model's prefill after a model switch. It runs on Hugging Face
Transformers; integration with vLLM is the work planned after the PoC.

\subsection{Risks}

The PoC is planned around the risks most likely to stop the project:

\begin{enumerate}
  \item \textbf{Translation quality.} A translated cache may not keep
  enough of the conversation for the target model to answer as it would
  after native prefill.
  \item \textbf{Training data.} Training the translator needs enough paired
  KV caches, produced by both models on the same text.
  \item \textbf{Net speed.} Translating, checking and loading a cache may
  take longer than native prefill, which would remove the benefit.
\end{enumerate}

\subsection{What the Demonstration Will Show}

The source model is SmolLM2-360M-Instruct and the target model is
Qwen2.5-0.5B-Instruct.

\begin{itemize}
  \item \textbf{Risk 1: translation quality.} Within a conversation, the
  router switches models and the target model continues from the
  translated cache. For the same prompts, it is compared with the target
  model after native prefill by first-token agreement, the length of the
  matching generated prefix, and answer F1 score.
  \item \textbf{Risk 2: training data.} The paired KV dataset built from
  HotpotQA passages and questions, and the trained SmolLM to Qwen
  translator with its development results, which are reported in the ML
  model report.
  \item \textbf{Risk 3: net speed.} TTFT against context length for the
  translated path, including translation and checking time, against native
  prefill, with peak GPU memory for both.
  \item Unsupported or rejected caches fall back to native prefill.
\end{itemize}

\subsection{If the PoC Fails}

If translated caches do not stay close to native prefill, or are not faster
than it, the team will narrow the project to reusing a model's own cache
across turns and measuring the cost of model switches. This stays in the
same problem area, and the router, runner and benchmark work carries over,
so work done before the PoC is not wasted. The exact rescoping would be
decided with the TA after the PoC demonstration.

\section{Expected Technology}

Table~\ref{TblTech} lists the tools we expect to use and why each was
chosen.

\begin{table}[H]
\caption{Expected technology} \label{TblTech}
\begin{tabularx}{\textwidth}{lX}
\toprule
\textbf{Need} & \textbf{Choice and reason}\\
\midrule
Language & Python for almost everything; C++ and CUDA (NVIDIA's GPU programming
platform) only if a custom kernel
turns out to be necessary.\\
Model framework & PyTorch and Hugging Face Transformers: the standard
tools for loading open models and inspecting their internals, which the
PoC needs.\\
Serving engine & Hugging Face Transformers for the PoC, then
vLLM~\cite{Kwon2023vLLM}, integrated rather than rewritten. vLLM has the
largest community and documentation. The cost is a large codebase to learn;
SGLang~\cite{Zheng2024SGLang} is the alternative if vLLM's design gets in
the way.\\
Models & SmolLM2-360M-Instruct and Qwen2.5-0.5B-Instruct, existing
open-weight models that are not retrained. The translator between them is
trained by the team.\\
Linter / formatter & \texttt{ruff} (lint and format), run in CI.\\
Unit tests & \texttt{pytest}, for the router, classifier, runner and
translator.\\
Code coverage & \texttt{coverage.py} through \texttt{pytest-cov}, reported
in CI; no target until the codebase exists.\\
CI & GitHub Actions for documents (running) and code (planned).\\
Performance tools & PyTorch Profiler and NVIDIA Nsight Systems, on the
department clusters (\texttt{ada}, 12$\times$ L40S, for development;
\texttt{grace}, 4$\times$ H100, for final benchmarks).\\
Editor & VS Code and Cursor, both locally and on the department clusters
through Remote-SSH, so the same editor is used on a laptop and on a GPU
machine.\\
Notebooks & Jupyter notebooks for quick demos and exploration. They run
locally, on Google Colab, or on a cluster through VS Code. Code that the
team keeps is moved into regular Python files with tests.\\
Google Colab & Quick GPU checks with its free GPU, for example while a
cluster is busy or not yet set up. Colab is not used for reported
measurements, because its GPU type, memory, and session length vary.\\
Project tools & git, GitHub, GitHub Projects, Microsoft Teams, Google
Meet.\\
\bottomrule
\end{tabularx}
\end{table}

\section{Use of AI and LLMs}

\begin{itemize}
  \item \textbf{Tools:} Claude Code and OpenAI Codex for code, debugging,
  and document drafting; GitHub Copilot code review on pull requests;
  Granola for meeting notes and summaries.
  \item \textbf{Code:} AI-written code is held to the same standard as any
  other code: it is reviewed in a PR and must pass the tests. For
  correctness-critical code, such as the cache translator, the author must
  be able to explain every line in review, since a plausible-looking
  result is not the same as a correct one.
  \item \textbf{Documents:} AI may draft or restructure text, but decisions
  (roles, schedule, scope) are made by the team, and every claim about
  other systems or papers is checked against the original source before
  it is submitted.
  \item \textbf{Disclosure:} each course document has an appendix
  describing how AI was used for it.
\end{itemize}

\section{Coding Standard}
\label{sec:standard}

Python code follows the PEP~8 style guide, enforced by \texttt{ruff} in CI.
Public functions
have type hints and a docstring. Tests live in \texttt{test/} and are named
after the module they test. If C++/CUDA code becomes necessary, it will
follow the Google C++ Style Guide, formatted with \texttt{clang-format}.

\bibliographystyle{plain}
\bibliography{../../refs/References}

\newpage{}

\section{Appendix --- Generative AI Use Disclosure}

This appendix describes how AI tools were used to write this document.
Planned use of AI during the project itself is in the Use of AI and LLMs
section.

\subsection{AI Tools Used}

\begin{itemize}
  \item Claude, through the Claude Code command-line tool (version
  2.1.283), using the Claude Sonnet 5 and Claude Opus 5.5 models for the
  draft and its revisions.
  \item GitHub Copilot pull request review, on PR \#3 of the project
  repository.
  \item Gemini 8 Flash, through Google AI Studio.
  \item For the November revision, Claude Code (version 2.1.329) with the
  Claude Opus 5.5 model.
\end{itemize}

\subsection{How and Where Used}

Table~\ref{TblAIUse} lists each part of the document that AI tools
contributed to, and how.

\begin{table}[H]
\caption{AI use by part of the document} \label{TblAIUse}
\begin{tabularx}{\textwidth}{lX}
\toprule
\textbf{Part of document} & \textbf{Use}\\
\midrule
Background research & Gemini was used to research the problem
area and find related papers on LLM serving and shared-prefix inference,
which informed the project framing and the related work named in this
plan.\\
All sections & Claude produced the first draft of every section from
information and decisions given by the team, and reorganized the document
to follow the course template.\\
Checklist and rubric & Claude compared the draft against
\path{docs/Checklists/DevPlan-Checklist.tex} and the course rubric and
listed the gaps (missing PoC fallback, technical and QA roles, schedule
timing, tool justification); the team decided how to fill each one.\\
PoC plan & Claude helped design the first PoC plan. For the November
revision, Claude rewrote Section~\ref{sec:poc}, the terms, technical
areas, schedule and technology for cross-model KV cache translation, from
the team's list of changes and decisions.\\
Schedule & Claude collected the course dates from the capstone
repository's weekly announcements and course schedule.\\
License & Copilot review pointed out that \texttt{LICENSE} was still the
template; the team replaced it.\\
Technology & Claude drafted the technology table and the bibliography
entries.\\
Reflection & The team supplied its own points and answers; Claude wrote
them up as prose, and the team checked that the result says what it
meant.\\
\bottomrule
\end{tabularx}
\end{table}

\subsection{Verification of AI-Generated Content}

William Pol drafted this document with Claude Code from decisions made by
the team, and Zayed Umer reviewed the pull request, requesting the PoC
memory/throughput note and the tooling additions reflected in this
version. The whole team reviews the document again before submission.

\begin{itemize}
  \item \textbf{Decisions:} meeting format, communication tool, role
  rotation, and fallback scope were chosen by the team; AI-proposed
  defaults were replaced or confirmed by the team before submission.
  \item \textbf{Results:} any correctness or performance numbers reported
  in future deliverables will come from running the code in the
  repository, not from the AI's description of it.
  \item \textbf{Dates:} course dates were checked against the capstone
  repository's announcements.
  \item \textbf{Bibliography:} the entries for the cited papers were
  drafted by the AI from memory, so each is checked against the paper
  itself before submission.
  \item \textbf{Related work:} the papers suggested by the AI are treated
  as leads to read, not as facts; none of their claims are relied on in
  this document until a team member has read the paper.
  \item \textbf{Review:} the full document was reviewed by team members
  against the checklist before submission.
\end{itemize}

The team is responsible for the accuracy, completeness, consistency, and
appropriateness of this document, whether material was written by a person
or an AI system.

\newpage{}

\section*{Appendix --- Reflection}

\input{../Reflection.text}

\begin{enumerate}
    \item \textbf{Why is it important to create a development plan prior to
    starting the project?}

    It is good for the team to review and agree on the terms of how
    we want to work together and develop this project, so that problems
    surface early instead of showing up in later deliverables. Making
    explicit decisions about roles, workflow, and technology, and
    committing to them as a team, is what prevents those problems in the
    first place: a decision that is only assumed rather than written down
    and agreed to is a decision the team can quietly disagree about later,
    right when there is no time left to resolve it. The PoC plan pushed
    this further: it required us to assess risk and think honestly about
    the weaknesses in our own proposal before we had written any code.
    That exercise directly shaped how we are developing the project, since
    it pushed us to write and run code early, so that if our core idea was
    going to fail, it would fail early, not right before a deadline.
    \item \textbf{In your opinion, what are the advantages and
    disadvantages of using CI/CD?}

    CI/CD is good for making the development pipeline concrete,
    giving the team an adequate, repeatable framework to develop against.
    It catches build and documentation errors, and later test failures,
    before they reach \texttt{main}, and it automates tedious tasks like
    PDF generation. However, we found CI/CD can be both flaky and too
    strict at the same time. In the main capstone repository, the
    pipeline for our project description submission would not let us
    merge because we had marked ``supervisor agreed'' as false, even
    though the instructor, Spencer Smith, had told us directly it was fine
    to proceed without a supervisor for now. Separately, GitHub Pages
    deployment for our own documentation has worked inconsistently, some
    runs succeeding and others not, for no clear reason. So while CI/CD is
    a useful tool, it can become a bottleneck rather than a help if it is
    not kept reliable, which ties back to the setup and maintenance cost
    of running it at all.
    \item \textbf{What disagreements did your group have in this
    deliverable, if any, and how did you resolve them?}

    Two disagreements came up while writing this deliverable. The
    first was whether to assign specific names to the newly proposed team
    roles right now. Some teammates argued that naming people would give
    the team firmer direction moving forward and make responsibilities
    clear. Others disagreed, arguing that attaching labels before we had
    actually worked together was misguided without a proper gauge of each
    teammate's strengths and weaknesses. We resolved this by deferring the
    decision rather than forcing an answer either way. The second was
    about the premise of the project itself: preliminary work for the PoC
    revealed that one of our core assumptions was weaker than we had first
    imagined. We agreed the overall direction of the project was still
    intact, because we have strong passion for this area, but the
    disagreement made clear that we need to keep testing our assumptions
    rather than assume the original framing still holds.
\end{enumerate}

\newpage{}

\section*{Appendix --- Team Charter}

This charter follows the structure of the University of Portland Team
Charter template.

\subsection*{External Goals}

Build something substantial enough to be a genuine portfolio and interview
talking point in systems engineering, and aim for a strong final grade;
present at the Capstone EXPO if the results are compelling.

\subsection*{Attendance}

These rules cover team meetings, TA meetings, and demonstrations.

\subsubsection*{Expectations}

Team members are expected to attend all scheduled team meetings on time and
stay for the full meeting; if unable to attend, notify the team as far in
advance as possible.

\subsubsection*{Acceptable Excuse}

Illness, other course or exam conflicts, and emergencies are acceptable
excuses. Simply forgetting, or deprioritizing without advance notice, is
not.

\subsubsection*{In Case of Emergency}

Notify the team chat and the meeting chair as soon as possible; the chair
will share Granola's meeting summary with the absent member to bring them
up to speed on decisions and action items.

\subsection*{Accountability and Teamwork}

These expectations apply to every member, whatever their current role.

\subsubsection*{Quality}

Members are expected to come to meetings having completed their agreed-upon
work, or to flag blockers before the meeting rather than at it. Pull
requests should be self-reviewed (tests passing, documents updated) before
requesting review.

\subsubsection*{Attitude}

We adopt the
\href{https://www.contributor-covenant.org/}{Contributor Covenant} as our
code of conduct. Disagreements are expected to be resolved through
discussion and, if unresolved, a team vote (see Decision Making below).

\subsubsection*{Stay on Track}

We check progress through commits, issue updates and work shared at
meetings. If a member makes no meaningful progress on their assigned work
for two weeks, or misses two consecutive team meetings without an
acceptable excuse (see Attendance above), the team will speak with them at
the next meeting. We will agree on next steps and a date to check progress.
If the problem continues for another two weeks, we will raise it with the
TA.

If a member misses an agreed deadline, they should update the task's issue
and tell the team what remains and what is blocking them. At the next
meeting, we will agree on a new deadline and, if needed, split up or
reassign the work. If they miss the new deadline without agreeing on an
extension with the team, we will raise it with the TA.

\subsubsection*{Team Building}

The team eats together after milestone deliverables (Problem Statement,
PoC demo, Rev 0 demo, Final demo) to mark the milestone and decompress
before the next one starts. Team meetings open with a few minutes of
non-project conversation before the agenda starts, so meetings stay a
place the team wants to show up to, not just a status check.

\subsubsection*{Decision Making}

Default to consensus; if consensus is not reached after discussion, decide
by majority vote among team members.

\end{document}
